\chapter{Common Errors and Their Remedies}\label{app:errors}
\markright{Common Errors and Their Remedies}

An error message is your friend, but at the beginning, reading one takes a
little practice. An error message that Salam gives before the program runs
usually has three parts: the message itself; the place of the error (the
file name, the line number and the column number; for example
\code{main.salam:3:5} means line 3, column 5); and a piece of that line with
a \code{\textasciicircum} mark under the spot where the problem was seen.
Most messages also carry a number, which appears in square brackets after
the word ``error''; for example \code{error[E001]}. Some messages end with a
line beginning \code{= help:} which makes a specific suggestion for fixing
the error. Any message that begins with \code{warning} is not an error: the
program runs, but Salam is telling you that something in it is probably
wrong. Runtime errors, such as running out of time, happen while the program
is running and have no \code{\textasciicircum} mark. This
appendix lists the most common messages that beginners meet, with their
cause and remedy; for brevity we have left out the piece of the program and
the \code{\textasciicircum}.

\section*{Unterminated string literal}

\begin{plainout}
error: unterminated string literal
\end{plainout}

\textbf{Cause.} The closing quotation mark of a string has been forgotten.
\textbf{Remedy.} On the line the message names, look at every string and
make sure it is opened and closed with \code{"}.

\section*{Expected \code{'end'} to close block}

\begin{plainout}
error: expected 'end' to close block (at end of file)
\end{plainout}

\textbf{Cause.} A block (a function, an \code{if}, a loop) has been opened
with a colon but its \code{end} has not been written. The line number
usually points at the end of the file, because Salam searched for the
\code{end} all the way down and did not find it.
\textbf{Remedy.} Count the colons that open blocks and count the
\code{end}s. Proper indentation makes the missing one easy to find.

\section*{Expected \code{':'} to open block}

\begin{plainout}
error: expected ':' to open block (at end of line)
\end{plainout}

\textbf{Cause.} The colon at the end of a \code{func main}, \code{repeat}
or \code{until} line has been forgotten. If the colon after the condition of
an \code{if} or \code{else} is missing, the message is slightly different:

\begin{plainout}
error: expected ':' after if condition (at end of line)
\end{plainout}

\textbf{Remedy.} Put the colon at the end of that line.

\section*{Unknown identifier}

\begin{plainout}
error[E001]: unknown identifier 'agee'
\end{plainout}

\textbf{Cause.} A name has been used that has not been defined. Most of the
time it is a spelling mistake (\code{agee} for \code{age}), or a variable
that was declared inside another block and is not available here.
\textbf{Remedy.} Compare the spelling of the name with its declaration. If
the variable is declared inside another loop or function, move it to the
right place.

\section*{Routine is not called}

\begin{plainout}
error[E110]: routine 'separator line' is not called: write 'separator line()' to call it
\end{plainout}

\textbf{Cause.} The name of a function has been written without
parentheses.
\textbf{Remedy.} Put parentheses after the name of the function:
\mbox{\code{separator line()}}.

\section*{Expected \code{','} between the values to print}

\begin{plainout}
error[E110]: expected ',' between the values to print
\end{plainout}

\textbf{Cause.} In a \code{println} or \code{print}, the comma between two
values has been forgotten; for example in the instruction
\mbox{\code{println "Age:" age}}.
\textbf{Remedy.} Put a comma between every two values:
\mbox{\code{println "Age:", age}}. The same error number also appears for
any instruction that does nothing; for example a line that is just
\code{a + 1} and puts its result nowhere.

\section*{A name cannot start with a digit}

\begin{plainout}
error: a name cannot start with a digit; ...
\end{plainout}

\textbf{Cause.} A variable name such as \code{2people} begins with a digit.
\textbf{Remedy.} Put the digit after a letter: \code{people2}.

\section*{Unused variable}

\begin{plainout}
error[E059]: unused variable 'age': ...
\end{plainout}

\textbf{Cause.} A variable has been declared but is never read anywhere.
\textbf{Remedy.} Either use it, or delete the line that declares it. If you
are keeping it only for a temporary experiment, you can start its name with
\code{\_} so that Salam passes over it; but as the message itself says, this
is not advised, and you should delete the line before you finish.

\section*{Unused function}

\begin{plainout}
error[E066]: unused function 'separator line': it is private and nothing calls it; ...
\end{plainout}

\textbf{Cause.} A function has been defined but is never called anywhere.
\textbf{Remedy.} Call it somewhere, or delete it. Do not forget that a call
needs parentheses: \mbox{\code{separator line()}}.

\section*{Cannot reassign immutable variable}

\begin{plainout}
error[E013]: cannot reassign immutable variable 'age'; ...
\end{plainout}

\textbf{Cause.} You have given a new value with \code{=} to a variable that
was declared without \code{mut}.
\textbf{Remedy.} On the declaration line, write \code{mut} before the name.
If the variable was really never meant to change, look again at the line
that gives it the new value; the mistake may be there.

\section*{Reserved word}

\begin{plainout}
error: 'package' is a reserved word, so it cannot be used as a name
  = help: pick a different name, for example 'package_'
\end{plainout}

\textbf{Cause.} The name of a variable or a function, or one of its words,
is one of the words of the language itself (Appendix~\ref{app:keywords}).
\textbf{Remedy.} Choose another name, or join the two words together.

\section*{Cannot assign type A to type B}

\begin{plainout}
error[E002]: cannot assign 'str' to 'i32'
\end{plainout}

\textbf{Cause.} A value of one type has been put where a value of another
type belongs: a string in a numeric variable, a floating-point number in an
integer variable, an argument of the wrong type passed to a function. The
type names in the message are the full names: \code{str}, \code{i32} (that
is, \code{int}), \code{f64} and \code{bool}. If you pass an argument of the
wrong type to a function, you get a similar message numbered \code{E012}:
``argument 1: cannot pass \code{'str'} to \code{'i32'}''.
\textbf{Remedy.} Make the types agree: with \code{as}, with
\code{to\_int()}, or by correcting the value.

\section*{Condition must be bool}

\begin{plainout}
error[E021]: if condition must be bool, got 'i32'
\end{plainout}

\textbf{Cause.} In an \code{if}, an \code{until} or the \code{? :}
operator, something has been written that is not a boolean value; for
example a number.
\textbf{Remedy.} Turn the condition into a comparison: instead of
\code{if count:} write \code{if count > 0:}.

\section*{An assignment cannot be used as a value}

\begin{plainout}
error[E013]: an assignment cannot be used as a value;
to compare two values write '==' (or 'eq') instead of '='
\end{plainout}

\textbf{Cause.} In a condition, \code{=} has been written instead of
\code{==}; for example \code{if age = 20:}. \code{=} assigns a value and
asks no question.
\textbf{Remedy.} Write \code{if age == 20:} or \code{if age eq 20:}.

\section*{Function can fall off its end without returning}

\begin{plainout}
error[E071]: function 'absolute' can fall off its end without returning
a 'i32' value; add a final 'ret'
\end{plainout}

\textbf{Cause.} A function that has declared a return type does not reach a
\code{ret} on one of its paths (for example when none of its conditions
holds). Up to version 0.4.9 this was only a warning and the program still
ran; from 0.5.0 it is an error and the program does not build, because a
function that reaches the end of its body has no value to give back.
\textbf{Remedy.} Put a \code{ret} for the default case at the end of the
function.

\section*{The program runs but never finishes}

\begin{plainout}
runtime error: execution timed out (possible infinite loop)
\end{plainout}

\textbf{Cause.} An infinite loop: the condition of an \code{until} never
becomes false. The online editor stops a program that takes longer than
three seconds with this message.
\textbf{Remedy.} Look in the loop for the variable that should have changed
in every iteration and did not. If you declared the loop variable with
\code{mut} but never changed it anywhere, Salam tells you before the program
runs:

\begin{plainout}
error[E069]: variable 'counter' is declared 'mut' but never mutated; ...
\end{plainout}

\section*{The program runs but the output is wrong}

The hardest kind of fault is one that Salam cannot detect, because the
program is correct as far as the language is concerned and simply does
something other than what you wanted. A few pointers:
\begin{itemize}
  \item Put temporary \code{println} instructions in several places in the
        program and look at the values of the variables in the middle of
        the work. Delete them once you have found the fault.
  \item Be suspicious of integer division (\code{7 / 2} is \code{3}) and of
        the order of operations.
  \item Make sure you have written \code{==} in conditions, not \code{=}.
  \item In loops, check where the counter starts and where it stops;
        ``off by one'' mistakes are very common.
  \item Make sure that accumulating variables are declared outside the
        loop.
  \item Try the program with the smallest possible input (one element,
        zero, a negative number).
\end{itemize}
